\section{Summary}
\noindent Complete connectomes now exist for insect nervous systems, but whether a fixed wiring diagram can learn, and what
its memories consist of, is unknown. Here, we trained a spiking model of the whole male \textit{Drosophila} central nervous
system, \val{L01} neurons with fixed wiring, to plan moves and to remember a visual cue. Spikes write the cue into neuronal
adaptation, which holds it silently until the next input reads it out. Cells that hold the trace are enriched for transcripts
of a calcium-activated potassium channel in fly cell atlases. The learned readout recalls the trace only in the brain state of
learning: the same trace evokes a different response at another delay, and a disease-like loss of potassium current rewrites
the trace and changes how it is read. Learning in the new state is spared. Memory failure can arise from a change of state
rather than a loss of storage.
